% ===== Edit metadata first when creating a new note =====
\newcommand{\papertitle}{在此填写论文标题}
\newcommand{\paperauthors}{作者 1, 作者 2, 作者 3}
\newcommand{\papervenue}{会议 / 期刊名称}
\newcommand{\paperyear}{2026}
\newcommand{\paperdoi}{10.xxxx/xxxxxx}
\newcommand{\paperurl}{https://example.org/paper}
\newcommand{\readingdate}{\today}
\newcommand{\paperonesentence}{用一句话概括本文的问题、方法和结论。}
\newcommand{\paperkeywords}{关键词 1; 关键词 2; 关键词 3}

\ifdefined\paperindexmode
  \paperindexentry{\papertitle}{\paperauthors}{\readingdate}{\paperonesentence}{\paperkeywords}
\expandafter\endinput
\fi

\documentclass[UTF8,fontset=fandol,zihao=-4]{ctexart}

% Shared preamble for paper-reading notes.
% Compile with XeLaTeX + biber for GB/T 7714-2015 references.

\usepackage[a4paper,margin=2.5cm,headsep=0.8cm,footskip=1cm]{geometry}
\usepackage{amsmath,amssymb,mathtools,bm}
\usepackage{graphicx}
\usepackage{booktabs,tabularx,array,longtable,multirow}
\usepackage{enumitem}
\usepackage{float}
\usepackage{caption}
\usepackage{fancyhdr}
\usepackage{xcolor}
\usepackage{hyperref}
\usepackage{bookmark}
\usepackage{csquotes}
\usepackage[backend=biber,style=gb7714-2015,sorting=none]{biblatex}

\hypersetup{
  colorlinks=true,
  linkcolor=black,
  citecolor=blue!50!black,
  urlcolor=blue!60!black,
  pdfcreator={XeLaTeX},
  pdfsubject={Paper Reading Notes}
}

\ctexset{
  section = {
    name = {, .},
    number = \arabic{section},
    format = \Large\bfseries
  },
  subsection = {
    format = \large\bfseries
  },
  subsubsection = {
    format = \normalsize\bfseries
  }
}

\setlength{\parindent}{2em}
\setlength{\parskip}{0.35em}
\linespread{1.2}
\setlist[itemize]{leftmargin=2em,itemsep=0.3em,topsep=0.3em}
\setlist[enumerate]{leftmargin=2em,itemsep=0.3em,topsep=0.3em}
\captionsetup{font=small,labelfont=bf}

% Keep page numbers stable across all pages, including the title page.
\setlength{\headheight}{14pt}
\pagestyle{fancy}
\fancyhf{}
\fancyfoot[C]{\thepage}
\renewcommand{\headrulewidth}{0pt}
\renewcommand{\footrulewidth}{0pt}
\fancypagestyle{plain}{
  \fancyhf{}
  \fancyfoot[C]{\thepage}
  \renewcommand{\headrulewidth}{0pt}
  \renewcommand{\footrulewidth}{0pt}
}

\newcolumntype{Y}{>{\raggedright\arraybackslash}X}
\newcommand{\NoteField}[2]{\textbf{#1} & #2 \\}

\addbibresource{../bib/references.bib}

\hypersetup{
  pdftitle={\papertitle},
  pdfauthor={\paperauthors}
}

\title{\papertitle\\\large 论文阅读笔记}
\author{\paperauthors}
\date{阅读日期：\readingdate}

\begin{document}

\maketitle

\section{基本信息}

\RenderBasicInfoTable

\section{快速浏览}

\subsection{一句话摘要}
\paperonesentence

\subsection{研究问题}
说明作者试图回答的核心研究问题，最好同时写清楚问题设定与约束条件。

\subsection{核心贡献}
\begin{itemize}
  \item 贡献 1：写清楚“提出了什么”。
  \item 贡献 2：写清楚“相比已有方法改进了什么”。
  \item 贡献 3：写清楚“为什么值得关心”。
\end{itemize}

\subsection{适用任务}
\begin{itemize}
  \item 任务类型
  \item 数据条件
  \item 是否支持迁移 / 少样本 / 跨被试等
\end{itemize}

\section{Method}

\subsection{总体框架}
建议从模型流程图或算法主线出发，先写输入到输出的主干。

\subsection{输入输出}
说明输入模态、时间窗、通道设定、输出任务形式。

\subsection{关键模块}
\begin{itemize}
  \item 模块 1：作用与输入输出
  \item 模块 2：和 baseline 的区别
  \item 模块 3：是否引入额外先验
\end{itemize}

\subsection{训练目标 / 损失函数}
如果论文有显式训练目标，建议直接抄录并解释每一项：
\begin{equation}
\mathcal{L} = \lambda_1 \mathcal{L}_1 + \lambda_2 \mathcal{L}_2.
\end{equation}

\subsection{关键公式}
保留你后续最可能复用的公式、符号定义或复杂推导入口。

\subsection{与已有方法的差异}
这里建议用“已有方法假设 / 本文改动 / 带来的影响”的格式写。

\section{Experiment}

\subsection{数据集}
记录数据来源、规模、预处理、切分方式和潜在偏差。

\subsection{Baseline}
列出比较对象，并说明这些 baseline 为什么合理。

\subsection{指标}
记录准确率、F1、AUC、kappa、MAE 等评价指标及其适用前提。

\subsection{主要结果}
优先写与本文核心 claim 直接相关的结果，不要只抄整张表。

\subsection{消融 / 可解释性（如有）}
记录最能说明方法成立性的证据。

\section{Reflection}

\subsection{优点}
\begin{itemize}
  \item 方法或实验上的可信优点
  \item 论文表达上的优点
\end{itemize}

\subsection{局限}
\begin{itemize}
  \item 数据、设定、可扩展性或泛化方面的限制
  \item 你对实验公平性的疑问
\end{itemize}

\subsection{对我当前研究的启发}
建议直接写成能够迁移到你当前课题的判断，不要只写“有帮助”。

\subsection{可迁移到 EEG foundation model 的点}
从预训练目标、跨数据集建模、被试泛化、标签效率等角度写。

\subsection{还没看懂的问题}
把真正卡住你的地方列出来，便于二次精读。

\section{Reuse for Proposal}

\subsection{可用于开题报告 / 综述的表述草稿}
尽量使用接近正式文风的完整句子，后续可直接改写迁移。

\subsection{可引用的背景动机}
记录能支撑“为什么要做这件事”的表述。

\subsection{可引用的方法描述}
记录能支撑“别人怎么做、与你有什么关系”的表述。

\section{References}

正文中请使用 \verb|\cite{...}| 引用，例如：相关工作可参考 \cite{deepseekai2025r1}。

\printbibliography[heading=none]

\end{document}
